\documentclass[conference]{IEEEtran}

\usepackage[utf8]{inputenc}
\usepackage[T1]{fontenc}
\usepackage{graphicx}
\usepackage{amsmath}
\usepackage{booktabs}
\usepackage{hyperref}
\usepackage{enumitem}

\title{Customer Churn Prediction and Agentic Retention Strategy System}

\author{
\IEEEauthorblockN{Harshit Jain, Vansh Dagar, Anishka Khurana}
\IEEEauthorblockA{Introduction to AI/ML \\ 4th Semester}
}

\begin{document}

\maketitle

\begin{abstract}
Customer churn prediction is a critical problem in customer relationship management, especially in the telecom industry. Traditional machine learning models can identify high-risk customers but fail to provide actionable insights for retention. This project presents an intelligent customer churn prediction system enhanced with an agentic AI layer that generates structured retention strategies. The system integrates supervised learning, retrieval-augmented generation (RAG), and LangGraph-based agent workflows. Experimental results demonstrate strong predictive performance with a Gradient Boosting model achieving a test ROC-AUC of 0.8419. The agentic system further transforms predictions into actionable recommendations, enabling a shift from predictive analytics to decision support.
\end{abstract}

\section{Introduction}

Customer churn is a major challenge for businesses, particularly in subscription-based industries such as telecommunications. Identifying customers who are likely to leave is essential for reducing revenue loss.

Traditional machine learning approaches can predict churn risk using historical customer data. However, these models do not provide actionable recommendations to retain customers. This creates a gap between prediction and decision-making.

To address this limitation, this project proposes an agentic AI-based system that combines machine learning predictions with retrieval-augmented reasoning. The system not only predicts churn risk but also generates structured retention strategies using domain knowledge.

This approach transforms churn prediction from a passive analytical tool into an active decision-support system.

\section{Methodology}

\subsection{Machine Learning Pipeline}

The dataset used is the Telco Customer Churn dataset. Data preprocessing includes handling categorical variables, encoding, and feature alignment.

Three models were evaluated:
\begin{itemize}[leftmargin=*]
\item Random Forest
\item Gradient Boosting
\item Logistic Regression
\end{itemize}

To address class imbalance, SMOTE was applied within the training pipeline. Model selection was performed using Stratified K-Fold Cross-Validation with ROC-AUC as the evaluation metric.

The Gradient Boosting model achieved the best performance with a CV ROC-AUC of 0.8453 and a test ROC-AUC of 0.8419.

\subsection{Agentic AI Workflow}

The agentic system is implemented using LangGraph, which enables a structured and stateful execution pipeline.

The workflow is defined using a custom state containing:
\begin{itemize}
\item Customer data
\item Churn probability
\item Key contributing factors
\item Retrieved contextual strategies
\item Final structured report
\end{itemize}

The system is composed of the following nodes:

\begin{itemize}
\item \textbf{Analyze Factors}: Extracts top churn drivers using feature importance data
\item \textbf{Retrieve Context}: Queries the vector database using semantic search
\item \textbf{Generate Strategy}: Produces structured recommendations using an LLM
\end{itemize}

A conditional edge is implemented to skip the agentic pipeline for low-risk customers, significantly reducing computational overhead and improving efficiency.

This design demonstrates the integration of symbolic workflow control with generative AI capabilities.

\subsection{Retrieval-Augmented Generation}

A local knowledge base of customer retention strategies is used. The document is split into chunks of size 300 with overlap 50.

Embeddings are generated using the all-MiniLM-L6-v2 model and stored in a Chroma vector database. The system retrieves the top 2 relevant chunks for each query.

This ensures that the LLM generates grounded and context-aware recommendations.
\subsection{System Architecture and Workflow}

The proposed system follows a hybrid architecture combining machine learning and agentic AI components. The workflow is designed to transform raw customer data into actionable business insights.

The complete pipeline operates as follows:

\begin{enumerate}
\item User provides customer data via the Streamlit interface
\item Data undergoes preprocessing including handling missing values and encoding
\item The trained Gradient Boosting model predicts churn probability
\item A threshold-based routing mechanism determines further execution
\item For high-risk customers:
\begin{itemize}
\item Key churn factors are extracted using feature importance
\item A semantic query is generated and passed to the RAG system
\item Relevant retention strategies are retrieved from the knowledge base
\item The LLM generates a structured retention report
\end{itemize}
\item Final results are displayed in a user-friendly dashboard
\end{enumerate}

This modular design ensures scalability, interpretability, and efficient resource usage.

\section{Results and UI Demonstration}

The system was tested using multiple real customer profiles from the Telco dataset. The machine learning model successfully predicts churn probability, which is further processed by the agentic system for decision-making.

\subsection{Prediction Performance}

The Gradient Boosting model achieved the best performance:
\begin{itemize}[leftmargin=*]
\item Cross-Validation ROC-AUC: 0.8453
\item Test ROC-AUC: 0.8419
\end{itemize}

These results demonstrate strong predictive capability in identifying customers at risk of churn.

\subsection{User Interface and Predictions}

The Streamlit-based interface allows users to input or select customer data and visualize churn risk using an intuitive gauge chart.

Example outputs:
\begin{itemize}[leftmargin=*]
\item Customer A:
\begin{itemize}
\item Predicted Label: Stay
\item Risk Probability: 4.81\%
\item Risk Category: Low Risk
\end{itemize}

\item Customer B:
\begin{itemize}
\item Predicted Label: Stay
\item Risk Probability: 2.19\%
\item Risk Category: Low Risk
\end{itemize}

\item Customer C:
\begin{itemize}
\item Predicted Label: Stay
\item Risk Probability: 16.82\%
\item Risk Category: Low Risk
\end{itemize}
\end{itemize}

\subsection{System Behaviour}

For customers below the churn threshold of 30\%, the system intelligently avoids triggering the agentic workflow. This optimization ensures efficient resource usage and reduces unnecessary API calls.

For high-risk customers, the system activates the LangGraph pipeline to:
\begin{itemize}[leftmargin=*]
\item Identify key contributing factors
\item Retrieve relevant retention strategies
\item Generate structured, actionable recommendations
\end{itemize}

The system introduces a conditional execution mechanism, where expensive LLM-based reasoning is only triggered for high-risk customers, making the solution both scalable and cost-efficient.

\subsection{Robustness Features}

The system demonstrates strong robustness through:
\begin{itemize}[leftmargin=*]
\item Handling missing and noisy data via preprocessing and imputation
\item Displaying preprocessing warnings to the user
\item Maintaining consistent feature alignment with trained model inputs
\end{itemize}

This ensures reliable predictions even when input data is incomplete or imperfect.

\section{Conclusion}

This project demonstrates how traditional churn prediction can be enhanced using agentic AI techniques. By integrating machine learning, RAG, and LangGraph workflows, the system provides actionable insights rather than just predictions.

The proposed system enables organizations to proactively retain customers and improve decision-making processes. Future work can focus on expanding the knowledge base, integrating real-time data, and deploying the system in enterprise environments.

\end{document}